\chapter{Metodología}

%%%%%%%%%%%%%%%%%%%%%%%%%%%%%%%%%%%%%%%%%
\section{Identificación de fallas}

\par De acuerdo al Capítulo 2, sección 2.1, las fallas en correas transportadoras se encuentran enmarcadas de acuerdo a \cite{khanna2022types} y para la implementación de modelos de reconocimiento de objetos que identifiquen estaas fallas, se consideraron las siguientes fallas:

\begin{itemize}
    \item Agujero (Hole): pérdida completa de material a través del espesor de la banda, de tamaño apreciable, originada por eventos de daño severo (impactos altos, atrapamiento de material o elementos extraños).
    \item Daño por impacto (Impact Damage): deformaciones, cortes o abrasiones localizadas en la zona de carga, producidas por el impacto repetitivo o puntual del material transportado.
    \item Perforación (Puncture): daño puntual en el que un elemento punzante atraviesa la banda generando una abertura de menor extensión superficial que un agujero, usualmente de carácter incipiente.
    \item Desgarro (Tear): separación longitudinal o transversal de la banda, asociada a sobrecargas, material atascado o concentraciones de tensión, y que puede propagarse a lo largo de la correa.
    \item Abrasión (Wear): desgaste progresivo de la superficie de la correa debido a fricción, que se manifiesta como adelgazamiento del recubrimiento y patrones característicos de pérdida de material.
\end{itemize}

\par De esta forma, las fallas antes mencionadas se convierten en \textit{clases} para los modelos de detección de objetos.

%%%%%%%%%%%%%%%%%%%%%%%%%%%%%%%%%%%%%%%%%
\section{Base de datos y entorno virtual}


\subsection{Base de datos elaborada}
\par El entrenamiento de modelos basados en redes neuronales (convolucionales o no) requiere disponer de un conjunto de datos suficientemente representativo, anotado y particionado en subconjuntos mutuamente excluyentes:

\begin{itemize}
    \item \textbf{Entrenamiento (train)}: conjunto utilizado para ajustar los parámetros del modelo mediante el proceso de optimización. Suele concentrar entre un 60\% y un 80\% del total de muestras disponibles.
    \item \textbf{Validación (valid)}: conjunto empleado durante el entrenamiento para calcular métricas intermedias, seleccionar hiperparámetros y monitorizar el sobreajuste. Sus ejemplos no se usan para actualizar los pesos del modelo y, por lo general, representan una fracción menor que el conjunto de entrenamiento.
    \item \textbf{Prueba (test)}: conjunto reservado para la evaluación final e independiente del modelo, una vez completado el ajuste y la selección de hiperparámetros. Permite estimar el desempeño esperado en condiciones de uso real.
\end{itemize}


\subsection{Recopilación y limpieza}

\par Mediante el sitio web \textit{Roboflow}, así como plataformas de datos abiertos tales como \textit{Kaggle Datasets}, \textit{ImageNet}, \textit{Open Images Dataset} y otros repositorios académicos de libre acceso, se recopilaron bases de datos relacionadas con fallas en correas transportadoras y con entornos industriales de transporte de material a granel. Cada conjunto presentaba su propia organización de clases, resolución de imágenes y formato de anotación, por lo que fue necesario unificar las etiquetas al formato YOLO antes de su utilización en los modelos de detección. \\

\par Posteriormente, se aplicó un proceso de depuración y limpieza de datos. En primer lugar, se descartaron imágenes irrelevantes (sin presencia clara de correa transportadora o sin defecto observable) y aquellas con problemas severos de calidad, como baja resolución, desenfoque o ruido excesivo. En segundo lugar, se normalizaron las anotaciones corrigiendo etiquetas inconsistentes, reclasificando ejemplos cuya clase no coincidía con el defecto presente y ajustando las cajas delimitadoras cuando su posición o tamaño eran inadecuados. Finalmente, se homogeneizó la nomenclatura de clases y el formato de las etiquetas, garantizando una base de datos coherente y consistente para el entrenamiento y evaluación de los modelos. \\

\par En total, se logró confeccionar una base de datos de 2312 imágenes, todas con resolución de 640x640px y etiquetas normalizadas de acuerdo al formato que usa YOLO.


%%%%%%%%%%%%%%%%%%%%%%%%%%%%%%%%%%%%%%%%%%%%%%%%%%%%%

\subsection{Entorno virtual}

\par Con el fin de asegurar la reproducibilidad de los experimentos y aislar las dependencias del proyecto, se implementó un entorno virtual basado en \textit{Python} en la raíz del repositorio \textit{Implementation-of-Object-Recognition-Algorithms-for-Conveyor-Belt-Defect-Detection}. Todas las librerías requeridas se gestionan mediante el archivo de requerimientos (\textit{Environment/requirements.txt}), que agrupa paquetes para aprendizaje profundo, procesamiento de datos, visión por computador y visualización de resultados. \\

\par La creación y activación del entorno se realiza con las utilidades estándar de \textit{Python} y la instalación automatizada de dependencias a partir de \textbf{Environment/requirements.txt}. La ejecución de los modelos se efectúa sobre una GPU AMD Radeon compatible con la pila ROCm/DirectML, utilizando la distribución oficial de PyTorch para esta plataforma.




%%%%%%%%%%%%%%%%%%%%%%%%%%%%%%%%%%%%%%%%%%%%%%%
\section{Implementación de YOLOv5}

\subsection{Estructura de trabajo}

\par La implementación de YOLOv5 desde el modelo oficial de Ultralytics en el ambiente actual sigue la estructura de trabajo mostrada en el Código \ref{lst:estructura-yolov5}. El modelo oficial de YOLOv5 se guarda en una carpeta dentro de YOLO. Posteriormente, se configuran todos los archivos necesarios a modo de capa envolvente del modelo original para poder utilizarlo de manera compatible al entorno virtual. \\


\begin{lstlisting}[
    float,
    floatplacement=H,                               
    caption={Estructura de directorios para la implementación de YOLO.},
    label={lst:estructura-yolov5},
    basicstyle=\ttfamily\scriptsize,
    numbers=none,                            % sin numeracion de lineas
    columns=fullflexible,
    frame=single,
    xleftmargin=2.5cm,
    framexleftmargin=0.5cm,
    framexrightmargin=0.5cm,
    breaklines=true,
    breakatwhitespace=true
]
YOLO/
|-- configs/
|   |-- dataset.yaml        : Rutas y clases del dataset
|   |-- model_variants.yaml : Parametros de escalado (variantes)
|   |-- train.yaml          : Hiperparametros de entrenamiento
|   `-- valid.yaml          : Hiperparametros de validacion
|-- engine/
|   |-- bn2gn_patch.py      : Parche para normalizacion por lotes
|   |-- bootstrap_miopen.py : Interfaz de mitigacion de MIOpen
|   |-- Trainer.py          : Modulo de entrenamiento
|   |-- Validator.py        : Modulo de validacion interna/externa
|   `-- warnings.py         : Modulo de advertencias y errores menores
|-- metrics/                : Carpeta de almacenamiento de metricas
|-- runs/                   : Carpeta de almacenamiento de ejecuciones
|-- utility/
|   |-- check_dataset.py    : Revisa formato dataset para modelo a entrenar
|   |-- clean_logs_runs.py  : Script de limpieza de logs y runs
|   |-- clean_metrics.py    : Script de limpieza de metricas
|   |-- clean_weights.py    : Script de limpieza de pesos registrados
|   `-- metrics.py          : Script de procesamiento de datos y post-procesado
|-- weights/                : Carpeta de almacenamiento de pesos finales
|-- yolov5/                 : Carpeta con el modelo oficial YOLOv5
|-- train.py                : Script principal de entrenamiento
|-- valid.py                : Script principal de validacion
`-- test.py                 : Script principal de pruebas
\end{lstlisting}

\par El flujo principal de experimentación se organiza en torno a los scripts \textit{train.py}, \textit{valid.py} y \textit{test.py}, que actúan como puntos de entrada al motor implementado en \textit{engine/}. El script \textit{train.py} lee la configuración desde \textit{configs/train.yaml}, \textit{configs/dataset.yaml} y \textit{configs/model\_variants.yaml}, instancia la variante de modelo seleccionada a partir del código base contenido en \textit{yolov5/} y delega el proceso de entrenamiento al módulo \textit{Trainer.py}. Durante este proceso se aplica el parche de normalización \textit{bn2gn\_patch.py} y la inicialización de MIOpen mediante \textit{bootstrap\_miopen.py}, se registran pérdidas y métricas por época en \textit{metrics/} y \textit{runs/}, y se almacenan los pesos intermedios y finales en \textit{weights/}. \\

\par Una vez completado el entrenamiento, el script \textit{valid.py} utiliza el módulo \textit{Validator.py} junto con la configuración de \textit{configs/valid.yaml} para cargar los pesos seleccionados desde \textit{weights/} y evaluar el modelo sobre el conjunto de validación, generando métricas agregadas (por ejemplo, precisión, exhaustividad y mAP) y salidas gráficas que se guardan nuevamente en \textit{metrics/} y \textit{runs/}. Finalmente, el script \textit{test.py} permite aplicar el modelo entrenado sobre la partición de prueba independiente, obteniendo predicciones cualitativas y cuantitativas que sirven para verificar la capacidad de generalización del detector en condiciones próximas al uso operativo. Los scripts auxiliares contenidos en \textit{utility/} se emplean para comprobar la coherencia del conjunto de datos y mantener ordenados los registros de ejecución, métricas y pesos generados a lo largo de los experimentos.

\subsubsection{Hiperparámetros de YOLOv5}

\par La configuración del entrenamiento de YOLOv5 considera un conjunto de hiperparámetros que controlan tanto la optimización del modelo como las funciones de pérdida y las técnicas de aumentación de datos. La Tabla \ref{tab:3.4-hiperparametros-ssd} (Ver Capítulo 5.1) resume los hiperparámetros principales seleccionados en esta memoria, junto con su valor fijado y una breve descripción de su rol dentro del proceso de entrenamiento. Estos valores se basan en configuraciones por defecto y en recomendaciones de ajuste de hiperparámetros propuestas para la familia Ultralytics YOLO \cite{ultralyticsHyperparameterTuning}.



%%%%%%%%%%%%%%%%%%%%%%%%%%%%%%%%%%%%%%%%%%%%%%%%%%%%
%%%%%%%%%%%%%%%%%%%%%%%%%%%%%%%%%%%%%%%%%%%%%%%%%%
\section{Implementación de SSD}

\subsection{Estructura de trabajo}

\par La implementación del detector \textit{Single Shot MultiBox Detector} (SSD) se integró reutilizando la misma estructura de trabajo definida para YOLOv5, de modo de compartir el entorno virtual, la configuración de experimentos y el flujo de entrenamiento/validación. En este caso, el código base de la arquitectura se organiza en la carpeta \textit{ssd/}, mientras que los módulos de configuración, motor y utilidades se mantienen en las mismas ubicaciones lógicas dentro de \textit{YOLO/}, tal como se muestra en el Código~\ref{lst:estructura-ssd}. \\

\begin{lstlisting}[
    float,
    floatplacement=H,                               
    caption={Estructura de directorios para la implementacion de SSD.},
    label={lst:estructura-ssd},
    basicstyle=\ttfamily\scriptsize,
    numbers=none,
    columns=fullflexible,
    frame=single,
    xleftmargin=2.5cm,
    framexleftmargin=0.5cm,
    framexrightmargin=0.5cm,
    breaklines=true,
    breakatwhitespace=true
]
SSD/
|-- configs/
|   |-- dataset.yaml        : Rutas y clases del dataset
|   |-- model_variants.yaml : Parametros de escalado (variantes)
|   |-- train.yaml          : Hiperparametros de entrenamiento
|   `-- valid.yaml          : Hiperparametros de validacion
|-- engine/
|   |-- bn2gn_patch.py      : Parche para normalizacion por lotes
|   |-- bootstrap_miopen.py : Interfaz de mitigacion de MIOpen
|   |-- Trainer.py          : Modulo de entrenamiento
|   |-- Validator.py        : Modulo de validacion interna/externa
|   `-- warnings.py         : Modulo de advertencias y errores menores
|-- metrics/                : Carpeta de almacenamiento de metricas
|-- runs/                   : Carpeta de almacenamiento de ejecuciones
|-- utility/
|   |-- check_dataset.py    : Revisa formato dataset para modelo a entrenar
|   |-- clean_logs_runs.py  : Script de limpieza de logs y runs
|   |-- clean_metrics.py    : Script de limpieza de metricas
|   |-- clean_weights.py    : Script de limpieza de pesos registrados
|   `-- metrics.py          : Script de procesamiento de datos y post-procesado
|-- weights/                : Carpeta de almacenamiento de pesos finales
|-- ssd/                    : Carpeta con el modelo oficial de SSD
|-- train.py                : Script principal de entrenamiento
|-- valid.py                : Script principal de validacion
`-- test.py                 : Script principal de pruebas
\end{lstlisting}

\par En el caso de SSD, los scripts \textit{train.py}, \textit{valid.py} y \textit{test.py} reutilizan el mismo motor definido en \textit{engine/}, modificando únicamente la lógica de construcción del modelo para instanciar la arquitectura SSD a partir del codigo contenido en \textit{ssd/}. El archivo \textit{train.py} lee las configuraciones de \textit{configs/train.yaml}, \textit{configs/dataset.yaml} y \textit{configs/model\_variants.yaml}, crea el modelo SSD con las cajas por defecto y capas de predicción correspondientes, y delega el entrenamiento al módulo \textit{Trainer.py}, registrando pérdidas, métricas y pesos en las carpetas \textit{metrics/}, \textit{runs/} y \textit{weights/}. De forma análoga, \textit{valid.py} y \textit{test.py} aprovechan \textit{Validator.py} y la misma estructura de carpetas para evaluar el modelo sobre los conjuntos de validación y prueba, manteniendo un flujo de experimentación homogéneo entre SSD y YOLOv5.

\subsubsection{Hiperparámetros de SSD}

\par La configuración del entrenamiento de SSD considera un conjunto de hiperparámetros que controlan la optimización del detector, la definición de las cajas por defecto y los criterios de correspondencia entre predicciones y anotaciones reales. La Tabla \ref{tab:3.4-hiperparametros-ssd} (ver Capítulo 5.1) resume los hiperparámetros principales seleccionados en esta memoria, junto con su valor fijado y una breve descripción de su rol dentro del proceso de entrenamiento. \\

\par En particular, se incluyen parámetros asociados a la tasa de aprendizaje base y su programación, el momento y la regularización, así como a las escalas y razones de aspecto de las cajas por defecto, los umbrales de IoU para la asignación de ejemplos positivos/negativos y las estrategias de aumentación de datos. Estos valores se basan en la configuración propuesta en el trabajo original de \textit{Single Shot MultiBox Detector} \cite{liu2016ssd}, adaptada al dominio de fallas en correas transportadoras para permitir una comparación controlada con los experimentos realizados con YOLOv5.
